\section{模型评价与推广}

\subsection{模型优点}

\textbf{一、口径统一、来源可核验。} 全文对 46 份数据文件逐份审计并给出哈希，把数据分为真实、半合成、插值、估算四级，每一处结论都注明其证据属于哪一级。特别是本文通过残差结构分析识别出标注为``真实训练日志''的 B1 实为公式生成数据，并据此调整了泛化检验的策略——若不区分，会把``还原生成参数''误当成``发现真实规律''。

\textbf{二、结论与证据强度匹配。} 对可以精确求解的部分（问题三的最优配置、问题二的弹性）给出解析参照与机器精度级的核验；对不可辨识的部分（配比的 $\omega$、桥接斜率、最优规模的平坦带）明确报告不可辨识，用区间或情景代替点估计，不用数值优化器的返回值掩盖证据不足。

\textbf{三、把``数据本身是否可信''纳入建模流程。} 本文没有把来源分级当作事后的说明文字，而是让诊断结果实际改变建模路径：B1 的合成性使本文改用 B4/B5 做泛化检验；B2 的坐标平移使其退出绝对比较；B8 的外推点使其退出 $\kappa$ 的分阶段估计。这些``负面结果''都保留在论文中。

\textbf{四、结构性结论对参数扰动稳健。} 问题三的``质量投入三段式''在三个成本函数、五档上下文长度下全部复现；问题二的 $\kappa\approx1$ 在 B6、B7 两个独立数据集上分别得 1.052 与 1.076。这些是本文最可信的结论。

\subsection{模型缺点与局限}

\textbf{一、质量数据的来源替换。} A1--A3 的原始文件在仓库中为 Git LFS 指针且实体未上传，本文按《数据说明》给出的公开来源同源重建。重建的 A2 为 17\,438 条（原记载 17\,523 条，差 0.5\%），A1 为 50\,412 条（目标 51\,230 条，book 域受下载预算限制少 788 条）。虽然字段结构与域级统计高度一致，但\textbf{抽样方式与原发布版不可能完全相同}，相关的绝对数值应以本文重建口径为准。

\textbf{二、质量域映射存在代理。} 17 个配方域中有 11 个没有对应的质量信号域，其质量值由文本体裁代理给定。这使得 $Q(p)$ 的绝对值带有未量化的系统偏差。本文的对策是只把 $Q(p)$ 用于单调折算，不做域级因果解读。

\textbf{三、配比因子不可辨识。} $h(p)$ 中的迁移强度 $\omega$ 缺少与附件 A 的共同实验，只能作为情景变量。这使问题三中最优配比 $p^\ast$ 无法内点求解，本文固定 $p$ 处理并在敏感性中覆盖 $\omega$ 的取值。

\textbf{四、桥接误差大。} Loss 与 Benchmark 之间的映射 $R^2$ 仅 0.16$\sim$0.25，留一 MAE 达 7.34 分。这直接限制了问题四预测的精度，也是 24 个月预测越过样本内上限却无法给出可信区间的根本原因。

\textbf{五、算力增长率估计不稳。} C1$\cap$C4 匹配子集时间跨度仅 0.76 年，本文改用 C4 全量历史估计，得 0.628 dex/年；但逐年最大值口径给出 0.480，两者差异反映估计的不稳定。本文以情景覆盖而非声称精确测得。

\textbf{六、外推占比高。} 问题三的全部最优配置都超出数据支持的 $N$、$D$ 范围，问题四的 24 个月预测超出样本内得分上限。这些结果的\textbf{结构性结论}（方向、次序、边界状态）比\textbf{数值本身}更可信。

\subsection{改进方向}

\textbf{一、补充质量信号的实测覆盖。} 若能获得 17 个配方域的实测质量评分，$Q(p)$ 的代理偏差即可消除，问题三中的配比内点优化也随之可行。

\textbf{二、设计配比与质量的联合实验。} 当前 $\omega$ 不可辨识的根源是附件 A 与 B 的实验相互独立。若能在同一实验中同时变动配比与质量，即可辨识 $h(p)$ 的迁移强度，把问题三的固定配比假设升级为联立优化。

\textbf{三、用更丰富的桥接数据降低映射误差。} 现有桥接的高可比样本仅 7 个且全部来自同一模型族。若能补充多族、多验证集的配对观测，桥接的 $R^2$ 有望显著提升，问题四的预测区间也随之收窄。

\textbf{四、建模参评模型数量效应。} 未来最高分取决于参评模型数，本文用 0.9 分位前沿回避了这个问题。若建立``提交数量$\to$极值分布''的模型，可以给出更贴近``最高分''的预测。

\textbf{五、引入真实的数据供给约束。} 本文按题面假设数据量无上限。现实中高质量数据存在供给上限，加入 $D\le D_{\max}$ 约束后，高预算区间的结构可能再次改变——这正是题面``数据墙''背景所指向的方向。

\subsection{推广应用}

本文的方法框架不限于语言模型，可推广到三类问题：

\textbf{一、算力与数据的采购决策。} 问题三的``质量投入启动预算''给出了一个可直接使用的判据：当预算低于某个阈值时，把钱花在清洗数据上不如扩大模型；超过阈值后质量投入的边际收益才转正。表~\ref{tab:q3_trans} 的启动点可直接作为决策参考。

\textbf{二、多源异构数据的融合评价。} ``组内平均$+$组间几何平均$+$凸组合消解''这一套评分框架，可直接迁移到任何``多个方向不一的指标需要合成单一得分、且指标间存在冲突''的场景，如教学质量评估、供应商评价、环境质量指数构建。

\textbf{三、技术进步的贡献分解。} 问题四``固定组成子集$+$对数线性模型$+$顺序无关分解''的做法，可推广到任何``规模扩张与技术改进共同驱动指标增长''的领域，如光伏组件效率、芯片制程算力、农作物单产。本文特别提示：\textbf{分解前必须控制组成变化}，否则会得出符号相反的结论。

\subsection{AI 使用披露}

按题面要求，本文披露人工智能工具的使用情况：

\begin{itemize}[leftmargin=2em,itemsep=2pt]
\item \textbf{所用工具}：Claude（Anthropic 公司），通过 MathModel 桌面端调用，模型为 Claude 系列。
\item \textbf{使用环节}：（1）代码调试与语法纠错；（2）绘图脚本的模板改写；（3）文字表达的润色与排版。
\item \textbf{未使用环节}：模型的数学形式选择、公式推导、参数辨识策略、结构性转移的判据设计、以及全部结论的取舍，均由参赛队完成。
\item \textbf{核验方式}：所有数值结果均由参赛队独立运行代码复算，解析参照、有限差分与资源守恒三类检验（见表~\ref{tab:verify}）用于交叉验证；AI 未参与任何结论性判断。
\item \textbf{文献}：全部参考文献均通过 Crossref 或 arXiv 官方接口核对元数据后引用，未使用 AI 生成的文献条目。
\end{itemize}
